\subsection{A forest obstruction from a completed short map}

\begin{proposition}[Proved: a forest in the short composition map]
\label{A7:short-map-forest}
Let $\mathcal C$ be a finite acyclic category with full support
$0<1<2<3<4$, and write $h_{ij}=|\mathcal C(i,j)|$.
Suppose $h_{01}\leq2$ and $h_{03}\leq2$.
On the full vertex set $U=\mathcal C(0,2)$, form the simple
undirected graph $G_P$ whose edges are the two-element image
sets of the maps
\[
 \mathcal C(0,1)\longrightarrow U,\qquad a\longmapsto b\circ a,
 \qquad b\in\mathcal C(1,2).
\]
Repeated unordered pairs give one edge, and a one-element image
gives no edge. If $G_P$ is a forest, every rational degree-two
cohomology product of $B\mathcal C$ is zero.
There is no restriction on $h_{12}$, $h_{02}$ or $h_{04}$.

Dually, suppose $h_{34}\leq2$ and $h_{14}\leq2$.
On the full vertex set $W=\mathcal C(2,4)$, form the simple
graph $G_R$ from the two-element image sets of
\[
 \mathcal C(3,4)\longrightarrow W,\qquad d\longmapsto d\circ c,
 \qquad c\in\mathcal C(2,3).
\]
If $G_R$ is a forest, the same product-vanishing conclusion holds.
\end{proposition}
\begin{proof}
Collapse $\mathcal C(0,4)$ to a singleton. By
Proposition~\ref{A7:extreme-arrow-collapse}, the quotient map
induces a surjection on second rational cohomology and preserves
products. The short map defining $G_P$ is unchanged. It therefore
suffices to prove product vanishing in this quotient.

The quotient is the left-act cone of its full subcategory on
$1,2,3,4$, with source stalks
\[
 A=\mathcal C(0,1),\quad U=\mathcal C(0,2),\quad
 L=\mathcal C(0,3),\quad \{*\}.
\]
Give $A$ and $L$ their unique trees with at most two vertices.
Extend the forest $G_P$, including all isolated vertices of $U$,
to a tree on exactly $U$. Such an extension exists: connect its
distinct connected components by arbitrary additional edges.
Every act map $A\to U$ sends its only possible edge to an edge
of this chosen tree or to a vertex, by the definition of $G_P$.
Every other nonidentity act map has target $L$ or $\{*\}$, so
it also sends each tree edge to an edge or a vertex.
The compatible-tree lemma, Lemma~\ref{A7:tree-act}, annihilates
all positive rational cohomology products of this act cone.
Pullback proves the asserted degree-two vanishing for
$\mathcal C$. Applying this argument to the opposite category
gives the sink statement.
\end{proof}

\begin{corollary}[Proved: an entire cylinder of short-map completions]
\label{A7:short-map-forest-prefix}
Fix a positive hom-set profile and a partial assignment of the
three short maps
\[
 P:A\times B\to U,\qquad Q:B\times C\to V,
 \qquad R:C\times D\to W.
\]
If every value of $P$ is assigned and its graph $G_P$ satisfies
the source hypothesis of Proposition~\ref{A7:short-map-forest},
then no completion of this partial assignment to an associative
category of that profile can model $\CP^2$.
The dual assertion holds when every value of $R$ is assigned and
its graph satisfies the sink hypothesis. In particular, the
source obstruction can be applied before any value of $Q$ or $R$
has been assigned.
\end{corollary}
\begin{proof}
The graph in the proposition depends only on the completed map
$P$, respectively $R$, and on the fixed hom-set cardinalities.
Every completion has exactly the same graph and obeys the same
cardinality hypotheses. Its rational second products therefore
vanish, which is incompatible with the projective-plane ring.
\end{proof}

\begin{remark*}
Completeness of the relevant short map is essential to this
certificate. A forest made only from some assigned columns can
acquire a cycle when another column is assigned. The test used
here first verifies that every entry of $P$, respectively $R$,
is known. It then retains the full vertex set, discards loops,
and deduplicates unordered edges before testing for a cycle.
The resulting obstruction removes every literal completion of
the recorded global prefix. Its unrestricted completion count
is the product of the appropriate hom-set sizes over all unknown
short-map entries; this number is not a count of associative
categories or of first-use-normalized representatives.
\end{remark*}
